\documentclass[twocolumn]{iopjournal}

\usepackage{amsmath,amssymb,amsfonts}
\usepackage{amsthm}

\fancyhfoffset[L]{0mm}
\setlength{\marginparwidth}{0pt}
\setlength{\marginparsep}{0pt}

% ==========================================
% ZENODO PREPRINT HEADERS (ACTIVE)
% ==========================================
\renewcommand{\articletype}[1]{%
  {\vspace*{-8mm}\noindent \Large \sf Preprint}\par
  \vspace*{6mm}\noindent{\scriptsize \sf{\bfseries \MakeUppercase{#1}}}}

\fancyhead[L]{}
\fancyhead[R]{}

% Theorem environments
\newtheorem{lemma}{Lemma}
\newtheorem{theorem}[lemma]{Theorem}
\theoremstyle{definition}
\newtheorem{definition}[lemma]{Definition}

\begin{document}

\articletype{Supplementary Material}

\title{The Gram Matrix of the Tetrad Equation}

\author{Stephen Nelson}

\affil{Field Dynamics Research}

\ead{stephen@fielddynamics.org}

\begin{abstract}
The tetrad equation $g_{\mu\nu} = \eta_{ab}\,e^a{}_\mu\,e^b{}_\nu$
defines the metric as a Gram matrix of the frame field.  This note
examines the Gram matrix as a mathematical object.  The equation is
even in the spatial frame: if $M_1$ is a solution, so is $-M_1$.  Two
solutions generate three Gram products --- self-coupling, self-coupling
of the inverted root, and cross-coupling --- whose properties follow
from matrix multiplication alone.  Reading the columns of the two
solutions as coordinates in $\mathbb{R}^3$ identifies the first with
a regular tetrahedron, the second with the inverted tetrahedron, and
their compound with a stellated octahedron.  Each $K \times K$ slice
of the affine connection decomposes canonically into $6 + 4 + 6 = 16$
entries; across $K = 4$ slices, the full connection decomposes as
$24 + 16 + 24 = 64$.  The torsion-free condition retains
$16 + 24 = 40$ and leaves 24 unretained.  The 40-component
torsion-free connection is not isomorphic to the full 64-component
connection: a tetrahedron is not a stellated octahedron.
Every step is a computation; every intermediate result is verifiable
with pencil and paper.
\end{abstract}

% ==================================================================

\section*{Premise}

The tetrad equation
\begin{equation}
  g_{\mu\nu} = \eta_{ab}\,e^a{}_\mu\,e^b{}_\nu
  \label{eq:tetrad}
\end{equation}
is quadratic in the frame field $e^a{}_\mu$: the metric depends on the
product of two copies of $e$.  In the spatial sector, the parity
$e^i{}_\mu \to -e^i{}_\mu$ sends both copies to their negatives; the
product is invariant.  The equation is therefore even in the spatial
frame: if $M_1$ is a solution, so is $M_2 = -M_1$.  The solution set
therefore contains both members of the parity pair; restricting to one
is a projection that requires independent justification.
The torsion-free condition
$\Gamma^\lambda_{\mu\nu} = \Gamma^\lambda_{\nu\mu}$, adopted in
general relativity in 1915,
provides that restriction at the connection level: it retains the
symmetric part of the lower index pair and discards the antisymmetric
part.  The vanishing of the antisymmetric components is a sufficient
condition for the Levi-Civita theorem~\cite{LeviCivita1917}, not a consequence of it.

This note examines what the tetrad equation contains before the
restriction is applied.  The method is to compute the Gram matrix that
the tetrad equation defines and to read its structure.  The results
follow from linear algebra.  No gravitational dynamics, field equations,
or physical assumptions beyond the tetrad equation itself are required.

% ==================================================================

\section*{Part I: The three Gram products}

This section is self-contained.  It requires only the tetrad equation
and matrix multiplication.

% ------------------------------------------------------------------

\subsection*{Definitions}

\begin{definition}[Gram matrix of the frame]
\label{def:gram}
Let $M$ denote the $D \times K$ matrix whose rows are the $D = 3$
spatial components of the $K = 4$ frame vectors.  The specific frame is
uniquely fixed by the Lorentz signature~\cite{Nelson2026a}.  The tetrad
equation~(\ref{eq:tetrad}) restricted to the spatial sector states
that the metric is the Gram matrix of the frame:
\begin{equation}
  G = M^T M.
  \label{eq:gram}
\end{equation}
$G$ is a $K \times K$ matrix.  Its entry
$G_{IJ} = \sum_{a=1}^{D} M_{aI}\,M_{aJ}$ is the inner product of frame
vector~$I$ with frame vector~$J$.
\end{definition}

\begin{definition}[Two solutions]
\label{def:roots}
Because Eq.~(\ref{eq:gram}) is quadratic in $M$, if $M_1$ is a
solution then $M_2 = -M_1$ is also a solution (spatial parity:
$e^i{}_\mu \to -e^i{}_\mu$, $i = 1,2,3$).  Both solutions satisfy the
same equation with the same metric.  (Time reversal acts on the temporal
leg, which is not included in the spatial matrix $M$; the two spatial
solutions form the complete parity pair.)  From two solutions, three
Gram products can be formed:
\begin{equation}
  G_{M_1 M_1} = M_1^T M_1, \quad
  G_{M_2 M_2} = M_2^T M_2, \quad
  G_{M_1 M_2} = M_1^T M_2.
  \label{eq:three_products}
\end{equation}
The transposed product $G_{M_2 M_1} = M_2^T M_1 = (-M_1)^T M_1
= -M_1^T M_1 = G_{M_1 M_2}$, so there are exactly three distinct
products.
\end{definition}

% ------------------------------------------------------------------

\subsection*{Lemmas}

\begin{lemma}[Root-blindness]
\label{lem:root_blind}
$G_{M_1 M_1} = G_{M_2 M_2}$.  The self-coupling Gram matrix is
identical for both solutions.
\end{lemma}

\begin{proof}
\begin{equation}
  G_{M_2 M_2} = M_2^T M_2 = (-M_1)^T(-M_1)
  = M_1^T M_1 = G_{M_1 M_1}. \qquad\square
  \label{eq:root_blind}
\end{equation}
\end{proof}

\noindent
The metric cannot distinguish the two roots.  It is a property of
matrix multiplication: the product $M^T M$ is even in $M$.  We call a
quantity \textit{root-blind} when it takes the same value on both
solutions.

\begin{lemma}[Cross-coupling sign reversal]
\label{lem:cross}
$G_{M_1 M_2} = -G_{M_1 M_1}$.  The cross-coupling Gram matrix negates
every entry of the self-coupling.
\end{lemma}

\begin{proof}
\begin{equation}
  G_{M_1 M_2} = M_1^T M_2 = M_1^T(-M_1)
  = -M_1^T M_1 = -G_{M_1 M_1}. \qquad\square
  \label{eq:cross}
\end{equation}
\end{proof}

\noindent
The self-coupling is root-blind; the cross-coupling is
\textit{root-resolving}: it is the only Gram product whose value
distinguishes $M_1$ from $M_2$.

\begin{lemma}[Explicit Gram matrix]
\label{lem:explicit}
For the regular tetrahedral frame in $D = 3$, the Gram matrix has
diagonal entries $D$ and off-diagonal entries $-1$:
\begin{equation}
  G_{IJ} = \begin{cases}
    D & \text{if } I = J, \\
    -1 & \text{if } I \neq J.
  \end{cases}
  \label{eq:explicit_G}
\end{equation}
\end{lemma}

\begin{proof}
The frame vectors uniquely fixed by the Lorentz
signature~\cite{Nelson2026a} are the columns of
\begin{equation}
  M_1 = \begin{pmatrix}
    +1 & +1 & -1 & -1 \\
    +1 & -1 & +1 & -1 \\
    +1 & -1 & -1 & +1
  \end{pmatrix}.
  \label{eq:M1}
\end{equation}
Each column is a frame vector with components $\pm 1$ in the three
spatial directions.  Compute $G = M_1^T M_1$ entry by entry.
The diagonal: $G_{II} = \sum_{a=1}^{3}(\pm 1)^2 = 3 = D$.  The
off-diagonal: for columns~1 and~2, the three component products are
$(+1)(+1) + (+1)(-1) + (+1)(-1) = 1 - 1 - 1 = -1$.  For columns~1
and~3: $(+1)(-1) + (+1)(+1) + (+1)(-1) = -1 + 1 - 1 = -1$.  In each
case, exactly one of the three spatial components agrees in sign and two
disagree, giving $G_{IJ} = -1$.  The same holds for all six off-diagonal
pairs, as can be verified column by column from~(\ref{eq:M1}).
\hfill$\square$
\end{proof}

\noindent
By Lemma~\ref{lem:root_blind}, $G_{M_2 M_2} = G_{M_1 M_1}$.  By
Lemma~\ref{lem:cross}, $G_{M_1 M_2} = -G_{M_1 M_1}$.  All three Gram
products are now determined:
\begin{equation}
  G_{M_1 M_1} = \begin{pmatrix}
    3 & -1 & -1 & -1 \\
    -1 & 3 & -1 & -1 \\
    -1 & -1 & 3 & -1 \\
    -1 & -1 & -1 & 3
  \end{pmatrix},\quad
  G_{M_2 M_2} = G_{M_1 M_1},\quad
  G_{M_1 M_2} = \begin{pmatrix}
    -3 & 1 & 1 & 1 \\
    1 & -3 & 1 & 1 \\
    1 & 1 & -3 & 1 \\
    1 & 1 & 1 & -3
  \end{pmatrix}.
  \label{eq:three_grams_explicit}
\end{equation}
Both $G_{M_1 M_1}$ and $G_{M_1 M_2}$ are symmetric matrices; the sign
of every entry is reversed.

\begin{lemma}[Orthogonal invariance]
\label{lem:orthogonal}
The three Gram products are invariant under $O(D)$.  If $M$ is replaced
by $RM$ for any orthogonal matrix $R$ ($R^T R = \mathbb{I}$), then
\begin{equation}
  (RM)^T(RM) = M^T R^T R\,M = M^T M.
  \label{eq:orthogonal_inv}
\end{equation}
The same identity holds for each of $G_{M_1 M_1}$, $G_{M_2 M_2}$, and
$G_{M_1 M_2}$: replacing $(M_1, M_2)$ with $(RM_1, RM_2)$ leaves all
three products unchanged.
\end{lemma}

\begin{proof}
$(RM_1)^T(RM_1) = M_1^T R^T R\,M_1 = M_1^T M_1$.
$(RM_2)^T(RM_2) = M_2^T R^T R\,M_2 = M_2^T M_2$.
$(RM_1)^T(RM_2) = M_1^T R^T R\,M_2 = M_1^T M_2$.
\hfill$\square$
\end{proof}

\noindent
The three Gram products are invariant under orthogonal changes of
representative and therefore constitute the complete
orthogonal-invariant matrix content of the parity pair $(M_1, M_2)$.

\begin{lemma}[Eigenstructure]
\label{lem:eigen}
The Gram matrix~(\ref{eq:three_grams_explicit}) has exactly two distinct
eigenvalues:
\begin{enumerate}
\item[(a)] Eigenvalue $0$ with multiplicity $1$ and eigenvector
  $(1,1,1,1)$.
\item[(b)] Eigenvalue $D + 1 = 4$ with multiplicity $D = 3$.
\end{enumerate}
\end{lemma}

\begin{proof}
(a)~Sum the entries in any row: $3 + (-1) + (-1) + (-1) = 0$.  Every
row sums to zero.  Therefore $G \cdot (1,1,1,1)^T = (0,0,0,0)^T$, and
$(1,1,1,1)$ is an eigenvector with eigenvalue~$0$.  The four frame
vectors sum to zero: $\sum_{I=1}^{K} M_{aI} = 0$ for each spatial
component~$a$.  This is closure: the columns of $M$ are linearly
dependent, with the all-ones vector spanning the null space.

(b)~Write $G = (D+1)\,\mathbb{I} - J$, where $\mathbb{I}$ is the
$K \times K$ identity and $J$ is the $K \times K$ all-ones matrix.  The
matrix $J$ has eigenvalue $K = 4$ with eigenvector $(1,1,1,1)$ and
eigenvalue $0$ with multiplicity $K - 1 = 3$.  Therefore, on the
$(K-1)$-dimensional subspace orthogonal to $(1,1,1,1)$:
$G = (D+1)\,\mathbb{I} - 0 = (D+1)\,\mathbb{I}$, giving eigenvalue
$D + 1 = 4$ with multiplicity~$3$.

Check: $\mathrm{tr}(G) = 4 \times 3 = 12$.  Sum of eigenvalues:
$0 + 3 \times 4 = 12$.  $\checkmark$ \hfill$\square$
\end{proof}

\noindent
The dimension $D = 3$ enters the construction of $M$ as the number of
spatial components; it reappears in the eigenvalue multiplicity as a
consistency check, confirming that the Gram matrix faithfully encodes the
spatial dimension of the frame.

% ==================================================================

\section*{Part II: Geometric identification}

By Lemma~\ref{lem:orthogonal}, the three Gram products depend only on
the orthogonal class of the chosen parity pair.  We therefore fix the
canonical representatives $M_1$ and $M_2 = -M_1$ and identify the
geometry they realize.

The canonical representatives are:
\begin{equation}
  M_1 = \begin{pmatrix}
    +1 & +1 & -1 & -1 \\
    +1 & -1 & +1 & -1 \\
    +1 & -1 & -1 & +1
  \end{pmatrix},\quad
  M_2 = \begin{pmatrix}
    -1 & -1 & +1 & +1 \\
    -1 & +1 & -1 & +1 \\
    -1 & +1 & +1 & -1
  \end{pmatrix}.
  \label{eq:M1M2}
\end{equation}
The inversion matrix
\begin{equation}
  I_{\mathrm{inv}} = \begin{pmatrix}
    -1 & 0 & 0 \\
    0 & -1 & 0 \\
    0 & 0 & -1
  \end{pmatrix}
  \label{eq:Iinv}
\end{equation}
satisfies $M_2 = I_{\mathrm{inv}}\,M_1 = -M_1$: multiplying every
entry of $M_1$ by $-1$ produces $M_2$.  This is verifiable by
inspection of~(\ref{eq:M1M2}).

The columns of $M_1$ and $M_2 = -M_1$ are disjoint: no column appears
in both matrices.  Together, the eight columns exhaust the $2^3 = 8$
sign combinations $(\pm 1,\,\pm 1,\,\pm 1)$ --- the complete vertex
set of a cube of side~2.  By Lemma~\ref{lem:explicit}, the four
columns of each matrix have constant pairwise inner product
$G_{IJ} = -1$ for all $I \neq J$.  Four vectors in
$\mathbb{R}^3$ with constant pairwise inner product form a regular
tetrahedron, unique up to rotation~\cite{Coxeter1973}.  The four
columns of $M_2 = I_{\mathrm{inv}}\,M_1 = -M_1$ form the inverted
tetrahedron.  This compound of two dual tetrahedra inscribed in a cube
is the stellated octahedron~\cite{Kepler1619,Cromwell1997}:
\begin{equation}
  S = \begin{bmatrix} M_1 & M_2 \end{bmatrix}.
  \label{eq:stellated}
\end{equation}

Following the coordinate-space identification established
in~\cite{Nelson2026a}, we denote the first tetrahedron as TetA (the
columns of $M_1$, the identity component under spatial parity) and the
second as TetB (the columns of $M_2 = -M_1$, the parity component).

% ==================================================================

\section*{Part III: The index grid}

The affine connection $\Gamma^\lambda_{\mu\nu}$ has three indices.
Fixing the upper index $\lambda$ yields a $K \times K$ matrix in the
lower pair $(\mu,\nu)$; we call this matrix a \textit{slice}.  There
are $K = 4$ slices, one for each value of $\lambda$.

Consider the $(\mu,\nu)$ index pair for a single slice.  Label the
four indices $\{t,x,y,z\}$.  The full $4 \times 4$ matrix of pairs is:
\begin{equation}
  \begin{array}{c|cccc}
    & t & x & y & z \\
    \hline
    t & tt & tx & ty & tz \\
    x & xt & xx & xy & xz \\
    y & yt & yx & yy & yz \\
    z & zt & zx & zy & zz
  \end{array}
  \label{eq:grid}
\end{equation}
Each $K \times K$ slice decomposes canonically into three sectors:
$K = 4$ diagonal entries, $K(K{-}1)/2 = 6$ strict upper-triangle
entries ($\mu < \nu$), and $K(K{-}1)/2 = 6$ strict lower-triangle
entries ($\mu > \nu$).  Total: $4 + 6 + 6 = 16 = K^2$.
Across $K = 4$ slices, the full affine connection decomposes as:
\begin{equation}
  \underbrace{4 \times 6}_{24}
  \;+\;
  \underbrace{4 \times 4}_{16}
  \;+\;
  \underbrace{4 \times 6}_{24}
  \;=\; 64.
  \label{eq:connection_decomp}
\end{equation}

The torsion-free condition
$\Gamma^\lambda_{\mu\nu} = \Gamma^\lambda_{\nu\mu}$ identifies each
lower-triangle entry with its upper-triangle mirror.  This retains the
16-component diagonal sector together with one 24-component off-diagonal
sector, giving $16 + 24 = 40$, and leaves the complementary
24-component sector unretained.

In the tetrad formulation, the connection $\Gamma^\lambda_{\mu\nu}$
is a function of the frame $M$, and its symmetric part, the
Christoffel connection, depends only on the metric $g = M^T M$,
which is root-blind by Lemma~\ref{lem:root_blind}.  The antisymmetric
part is the complement and carries the root-resolving content of
Lemma~\ref{lem:cross}.  The torsion-free projection produces the
40-component sector (16 shared plus 24 retained TetA) and discards
the complementary 24-component TetB sector.

% ==================================================================

\section*{Part IV: The non-isomorphism}

\begin{theorem}[Torsion-free projection discards TetB]
\label{thm:noniso}
In the tetrad formulation, the torsion-free condition
$\Gamma^\lambda_{\mu\nu} = \Gamma^\lambda_{\nu\mu}$ retains the
16-component shared sector together with the 24-component TetA
sector, and discards the complementary 24-component TetB sector.
The torsion-free connection is therefore a projection of the full
64-component two-root affine connection, not an equivalent
rewriting of it.
\end{theorem}

\begin{proof}
By Lemma~\ref{lem:root_blind}, the metric $g = M^T M$ is root-blind,
so the Christoffel sector built from $g$ is root-blind.  By
Part~III, the torsion-free projection retains exactly the diagonal
sector together with one strict triangular sector, giving
$16 + 24 = 40$ components.  By Lemma~\ref{lem:cross}, complementary
root-resolving content exists and is nonzero.  By the Part~II
identification of the two root branches as TetA and TetB, the
retained 24-component sector is TetA and the complementary
discarded 24-component sector is TetB.  Hence the torsion-free
construction is a projection of the full $24 + 16 + 24 = 64$
two-root connection:
\begin{equation}
  \boxed{16 + 24 = 40 \;\neq\; 64 = 24 + 16 + 24.}
  \label{eq:noniso}
\end{equation}
\hfill$\square$
\end{proof}

% ==================================================================

\section*{Summary}

\vspace{6pt}
\noindent
{\footnotesize
\begin{tabular*}{\columnwidth}{@{\extracolsep{\fill}}lll@{}}
\hline
Result & Proves & Requires \\
\hline
\multicolumn{3}{l}{\textit{Part I: The three Gram products}} \\
Lemma 1 & Root-blindness & Def.\ of $G = M^T M$ \\
Lemma 2 & Cross-coupling sign & Def.\ of $G_{M_1 M_2}$ \\
Lemma 3 & Explicit $G$ & $M$ from~\cite{Nelson2026a} \\
Lemma 4 & Orthogonal invariance & Def.\ of $G = M^T M$ \\
Lemma 5 & Eigenstructure & Lemma 3 \\[4pt]
\multicolumn{3}{l}{\textit{Part II: Geometric identification}} \\
Tetrahedron & $G_{IJ} = -1$ $\Rightarrow$ reg.\ tet. & Lemma 3, Coxeter \\
Inversion & $M_2 = I_{\mathrm{inv}}\,M_1$ & Def.\ 2 \\
Compound & $S = [M_1 \mid M_2]$ = stell.\ oct. & Kepler~\cite{Kepler1619} \\[4pt]
\multicolumn{3}{l}{\textit{Part III: The index grid}} \\
Grid & $6 + 4 + 6 = 16$ per slice & Matrix structure \\
Decomp. & $24 + 16 + 24 = 64$ & $\times\,K = 4$ slices \\
Torsion-free & $16 + 24 = 40$ retained & $\Gamma^\lambda_{\mu\nu} = \Gamma^\lambda_{\nu\mu}$ \\[4pt]
\multicolumn{3}{l}{\textit{Part IV: The non-isomorphism}} \\
Thm.\ 1 & Retains shared + TetA, drops TetB & Parts I--III \\
\hline
\end{tabular*}}
\vspace{6pt}

\noindent
The tetrad equation is even in the spatial frame: solutions come in
parity pairs.  Two roots produce three Gram products, displayed as
matrices in Part~I.  The three products are orthogonal-invariant
(Lemma~4): any choice of representative from each parity component
produces the same matrices.  Part~II chooses canonical representatives
and identifies the two roots with two tetrahedra whose compound is a
stellated octahedron.  Part~III decomposes the full connection as
$24 + 16 + 24 = 64$; the torsion-free condition retains $16 + 24 = 40$
and leaves 24 unretained.
Part~IV concludes that the torsion-free projection retains the shared
sector together with TetA and discards the complementary TetB sector.
Equivalently, $16 + 24 = 40 \neq 64 = 24 + 16 + 24$, so the
torsion-free construction does not preserve the full compound
structure defined by the tetrad equation.

% ==================================================================

\begin{thebibliography}{99}

\bibitem{Nelson2026a}
S.~Nelson,
\textit{The Coordinate-Space Geometry of the Tetrad Frame},
Field Dynamics Research (2026), submitted.

\bibitem{LeviCivita1917}
T.~Levi-Civita,
\textit{Nozione di parallelismo in una variet\`{a} qualunque e
conseguente specificazione geometrica della curvatura Riemanniana},
Rend.\ Circ.\ Mat.\ Palermo \textbf{42} (1917) 173.

\bibitem{Kepler1619}
J.~Kepler,
\textit{Harmonices Mundi} (1619).

\bibitem{Cromwell1997}
P.~R.~Cromwell,
\textit{Polyhedra}
(Cambridge University Press, Cambridge, 1997).

\bibitem{Coxeter1973}
H.~S.~M.~Coxeter,
\textit{Regular Polytopes},
3rd ed.\ (Dover, New York, 1973).

\end{thebibliography}

\end{document}
